\frfilename{mt247.tex} 
\versiondate{26.8.13} 
\copyrightdate{1995} 
      
\def\chaptername{Ruang fungsi} 
\def\sectionname{Kekompakan lemah dalam $L^1$} 
      
\newsection{247} 
      
Sekarang kita sampai pada ciri keterintegralan seragam yang paling mencolok:
sifat ini memberikan karakterisasi bagi subhimpunan $L^1$ yang relatif kompak
dalam topologi lemah (247C). Saya menempatkannya dalam bagian tersendiri karena
pembahasannya memerlukan sedikit pengetahuan tentang analisis fungsional --
khususnya, tentu saja, tentang topologi lemah pada ruang Banach. Saya akan
berusaha menyajikannya dengan istilah yang dapat dipahami oleh pembaca yang
baru mengenal teori ruang bernorma, sebab hasil ini pada hakikatnya merupakan
hasil teori ukur dan sekaligus sangat penting bagi penerapan dalam teori
peluang. Definisi-definisi pokoknya telah saya tuliskan dalam
\S\S2A3-2A5.  %2A3 2A4 2A5 
      
\leader{247A}{}\cmmnt{ Sebagian argumen untuk teorema utama di bawah akan
lebih lancar jika saya terlebih dahulu memisahkan sebuah gagasan yang pada
dasarnya merupakan kasus khusus sederhana dari tema yang berulang dalam
latihan-latihan bab ini (241Yg, 242Yb, 243Ya, 244Yd).
      
\medskip 
      
\noindent}{\bf Lema} Misalkan $(X,\Sigma,\mu)$ suatu ruang ukur dan $G$
sebarang anggota $\Sigma$. Misalkan $\mu_G$ ukuran subruang pada
$G$\cmmnt{, sehingga $\mu_GE=\mu E$ apabila $E\subseteq G$ dan
$E\in\Sigma$}. Tetapkan
      
\Centerline{$U=\{u:u\in L^1(\mu),\,u\times\chi G^{\ssbullet}=u\} 
\subseteq L^1(\mu)$.} 
      
\noindent Maka terdapat suatu isomorfisme $S$ antara ruang bernorma terurut
$U$ dan $L^1(\mu_G)$, yang diberikan oleh
      
\Centerline{$S(f^{\ssbullet})=(f\restr G)^{\ssbullet}$} 
      
\noindent untuk setiap $f\in\eusm L^1(\mu)$ dengan
$f^{\ssbullet}\in U$.
      
\proof{ Pertama-tama perlu dinyatakan secara eksplisit bahwa $U$ adalah
subruang linear dari $L^1(\mu)$. Integrasi pada subruang telah dibahas dalam
\S\S131 dan 214; khususnya, telah dicatat bahwa $f\restr G$ terintegralkan dan
bahwa
      
\Centerline{$\int|f\restr G|d\mu_G 
=\int|f|\times\chi G\,d\mu\le\int|f|d\mu$} 
      
\noindent untuk setiap $f\in\eusm L^1(\mu)$ (131Fa). Jika
$f$, $g\in\eusm L^1(\mu)$ dan $f=g\,\,\mu$-hampir di mana-mana, maka
$f\restr G=g\restr G\,\,\mu_G$-hampir di mana-mana; jadi rumus yang diajukan
untuk $S$ memang mendefinisikan suatu peta dari $U$ ke $L^1(\mu_G)$.
      
Karena
      
\Centerline{$(f+g)\restr G=(f\restr G)+(g\restr G)$, 
\quad$(cf)\restr G=c(f\restr G)$} 
      
\noindent untuk semua $f$, $g\in\eusm L^1(\mu)$ dan semua $c\in\Bbb R$,
$S$ linear. Karena
      
\Centerline{$f\le g\,\,\mu$-hampir di mana-mana$\,\Longrightarrow\, 
f\restr G\le g\restr G\,\,\mu_G$-hampir di mana-mana,} 
      
\noindent $S$ mempertahankan urutan. Karena
$\int|f\restr G|d\mu_G\le\int|f|d\mu$ untuk setiap
$f\in\eusm L^1(\mu)$, berlaku $\|Su\|_1\le\|u\|_1$ untuk setiap $u\in U$.
      
Untuk melihat bahwa $S$ surjektif, ambil sebarang $v\in L^1(\mu_G)$.
Nyatakan $v$ sebagai $g^{\ssbullet}$ dengan $g\in\eusm L^1(\mu_G)$.
Menurut 131E, $f\in\eusm L^1(\mu)$, dengan $f(x)=g(x)$ untuk
$x\in\dom g$ dan $0$ untuk $x\in X\setminus G$; akibatnya
$f^{\ssbullet}\in U$, $f\restr G=g$, dan
$v=S(f^{\ssbullet})\in S[U]$.
      
Untuk melihat bahwa $S$ mempertahankan norma, perhatikan bahwa untuk setiap
$f\in\eusm L^1(\mu)$,
      
\Centerline{$\int|f\restr G|d\mu_G 
=\int |f|\times\chi G\,d\mu$,} 
      
\noindent sehingga, jika $u=f^{\ssbullet}\in U$, diperoleh
      
\Centerline{$\|Su\|_1 
=\int|f\restr G|d\mu_G 
=\int|f|\times\chi G\,d\mu 
=\|u\times\chi G^{\ssbullet}\|_1 
=\|u\|_1$.} 
} 
      
\leader{247B}{Akibat} Misalkan $(X,\Sigma,\mu)$ sebarang ruang ukur dan
$G\in\Sigma$ suatu himpunan terukur yang dapat dinyatakan sebagai gabungan
terhitung dari himpunan-himpunan berukuran hingga. Definisikan $U$ seperti
dalam 247A, dan misalkan $h:L^1(\mu)\to\Bbb R$ sebarang fungsional linear
kontinu. Maka terdapat $v\in L^{\infty}(\mu)$ sedemikian sehingga
$h(u)=\int u\times v\,d\mu$ untuk setiap $u\in U$.
      
\proof{ Misalkan $S:U\to L^1(\mu_G)$ isomorfisme yang diuraikan dalam 247A.
Karena $S^{-1}:L^1(\mu_G)\to U$ linear dan kontinu, $h_1=hS^{-1}$ merupakan
anggota dual ruang bernorma $(L^1(\mu_G))^*$ dari $L^1(\mu_G)$. Ukuran $\mu_G$
tentu saja bersifat $\sigma$-hingga, dan karena itu dapat dilokalkan (211L),
sehingga menurut 243Gb terdapat $v_1\in L^{\infty}(\mu_G)$ sedemikian sehingga
      
\Centerline{$h_1(u)=\int u\times v_1d\mu_G$} 
      
\noindent untuk setiap $u\in L^1(\mu_G)$.
      
Nyatakan $v_1$ sebagai $g_1^{\ssbullet}$, dengan $g_1:G\to\Bbb R$ fungsi
terukur terbatas. Tetapkan $g(x)=g_1(x)$ untuk $x\in G$, dan $0$ untuk
$x\in X\setminus G$; maka $g:X\to\Bbb R$ adalah fungsi terukur terbatas dan
$v=g^{\ssbullet}\in L^{\infty}(\mu)$. Jika $u\in U$, nyatakan $u$ sebagai
$f^{\ssbullet}$ dengan $f\in\eusm L^1(\mu)$; maka
      
$$\eqalignno{h(u) 
&=h(S^{-1}Su) 
=h_1((f\restr G)^{\ssbullet}) 
=\int(f\restr G)\times g_1d\mu_G\cr 
&=\int(f\times g)\restr G\,d\mu_G 
=\int f\times g\times\chi G\,d\mu 
=\int f\times g\,d\mu 
=\int u\times v.\cr}$$ 
      
\noindent Karena $u$ sebarang, hasilnya terbukti.
}  
      
\leader{247C}{Teorema} Misalkan $(X,\Sigma,\mu)$ sebarang ruang ukur dan
$A$ suatu subhimpunan dari $L^1={L}^1(\mu)$. Maka $A$ dapat diintegralkan
secara seragam jika dan hanya jika himpunan tersebut relatif kompak dalam $L^1$ untuk
topologi lemah pada $L^1$.
      
\proof{{\bf (a)} Andaikan $A$ relatif kompak dalam topologi lemah. Saya akan
menunjukkan bahwa himpunan tersebut memenuhi syarat (iii) dalam 246G.
      
\medskip 
      
\quad{\bf (i)} Jika $F\in\Sigma$, tentu saja
$\sup_{u\in A}|\int_Fu|<\infty$, karena $u\mapsto\int_Fu$ termasuk dalam
$(L^1)^*$, dan untuk setiap $h\in (L^1)^*$ citra suatu himpunan yang relatif
kompak dalam topologi lemah di bawah $h$ harus terbatas (2A5Ie).
      
\medskip 
      
\quad{\bf (ii)} Sekarang andaikan $\sequencen{F_n}$ suatu barisan himpunan
yang saling lepas dalam $\Sigma$. \Quer\ Andaikan, jika mungkin, bahwa
\discrcenter{468pt}{$\sequencen{\sup_{u\in A}|\int_{F_n}u|}$ }tidak konvergen ke $0$. 
Maka terdapat suatu barisan naik tegas $\sequence{k}{n(k)}$ dalam $\Bbb N$
sedemikian sehingga
      
\Centerline{$\gamma 
=\Bover12\inf_{k\in\Bbb N}\sup_{u\in A}|\int_{F_{n(k)}}u|>0$.} 
      
\noindent Untuk setiap $k$, pilih $u_k\in A$ sedemikian sehingga
$|\int_{F_{n(k)}}u_k|\ge\gamma$. Karena $A$ relatif kompak dalam topologi
lemah, terdapat suatu titik klaster $u$ dari $\sequence{k}{u_{k}}$ dalam
$L^1$ untuk topologi lemah (2A3Ob). Tetapkan $\eta_j=2^{-j}\gamma/6>0$ untuk
setiap $j\in\Bbb N$.
      
Sekarang kita dapat memilih secara induktif suatu barisan naik tegas
$\sequence{j}{k(j)}$ sedemikian sehingga, untuk setiap $j$,
      
\Centerline{$\int_{F_{n(k(j))}}(|u|+\sum_{i=0}^{j-1}|u_{k(i)}|) 
\le\eta_j$} 
      
\Centerline{$\sum_{i=0}^{j-1} 
|\int_{F_{n(k(i))}}u-\int_{F_{n(k(i))}}u_{k(j)}|\le\eta_j$} 
      
\noindent untuk setiap $j$, dengan $\sum_{i=0}^{-1}$ ditafsirkan sebagai $0$.
\Prf\ Andaikan
$\langle k(i)\rangle_{i<j}$ telah dipilih, dan tetapkan
$v^*=|u|+\sum_{i=0}^{j-1}|u_{k(i)}|$. Maka
$\lim_{k\to\infty}\int_{F_{n(k)}}v^*=0$, misalnya menurut Teorema
Konvergensi Terdominasi Lebesgue. Karena itu terdapat $k^*$ sedemikian sehingga
$k^*>k(i)$ untuk setiap $i<j$ dan $\int_{F_{n(k)}}v^*\le\eta_j$ untuk setiap
$k\ge k^*$. Selanjutnya, peta
      
\Centerline{$w\mapsto\sum_{i=0}^{j-1} 
|\int_{F_{n(k(i))}}u-\int_{F_{n(k(i))}}w|:L^1\to\Bbb R$} 
      
\noindent kontinu dalam topologi lemah pada $L^1$ dan bernilai nol di $u$.
Karena $u$ merupakan titik klaster dari setiap ekor
$\{u_{k}:k\ge k^*\}$, terdapat $k(j)\ge k^*$ sedemikian sehingga
$\sum_{i=0}^{j-1}|\int_{F_{n(k(i))}}u-
\int_{F_{n(k(i))}}u_{k(j)}|<\eta_j$. Ini melanjutkan konstruksi. \Qed
      
Misalkan $v$ sebarang titik klaster dalam $L^1$, untuk topologi lemah, dari
$\sequence{j}{u_{k(j)}}$. Dengan menetapkan $G_i=F_{n(k(i))}$, kita mempunyai
$|\int_{G_i}u-\int_{G_i}u_{k(j)}|\le\eta_j$ apabila $i<j$. Jadi,
$\lim_{j\to\infty}\int_{G_i}u_{k(j)}$ ada $=\int_{G_i}u$ untuk setiap
$i$, serta $\int_{G_i}v=\int_{G_i}u$ untuk setiap $i$; dengan menetapkan
$G=\bigcup_{i\in\Bbb N}G_i$,
      
\Centerline{$\int_Gv=\sum_{i=0}^{\infty}\int_{G_i}v 
=\sum_{i=0}^{\infty}\int_{G_i}u=\int_Gu$,} 
      
\noindent berdasarkan 232D, karena $\sequence{i}{G_i}$ saling lepas.
      
Untuk setiap $j\in\Bbb N$,
      
$$\eqalignno{\sum_{i=0}^{j-1}|\int_{G_i}u_{k(j)}| 
&+\sum_{i=j+1}^{\infty}|\int_{G_i}u_{k(j)}|\cr 
&\le\sum_{i=0}^{j-1}\int_{G_i}|u| 
+\sum_{i=0}^{j-1}|\int_{G_i}u-\int_{G_i}u_{k(j)}| 
+\sum_{i=j+1}^{\infty}\int_{G_i}|u_{k(j)}|\cr 
&\le\sum_{i=0}^{j-1}\eta_i+\eta_j+\sum_{i=j+1}^{\infty}\eta_i 
=\sum_{i=0}^{\infty}\eta_i 
=\Bover{\gamma}{3}.\cr}$$ 
      
\noindent Di sisi lain, $|\int_{G_j}u_{k(j)}|\ge\gamma$. Karena itu
      
\Centerline{$|\int_Gu_{k(j)}| 
=|\sum_{i=0}^{\infty}\int_{G_i}u_{k(j)}| 
\ge\Bover23\gamma$.} 
      
Pernyataan ini berlaku untuk setiap $j$. Karena setiap himpunan terbuka lemah
yang memuat $v$ beririsan dengan $\{u_{k(j)}:j\in\Bbb N\}$, kita memperoleh
$|\int_Gv|\ge\bover23\gamma$ dan $|\int_Gu|\ge\bover23\gamma$. Akan tetapi,
      
\Centerline{$|\int_Gu| 
=|\sum_{i=0}^{\infty}\int_{G_i}u| 
\le\sum_{i=0}^{\infty}\int_{G_i}|u| 
\le\sum_{i=0}^{\infty}\eta_i 
=\Bover{\gamma}{3}$,} 
      
\noindent suatu kemustahilan. \Bang
      
Kontradiksi ini menunjukkan bahwa
$\lim_{n\to\infty}\sup_{u\in A}|\int_{F_n}u|=0$. Karena
$\sequencen{F_n}$ sebarang, $A$ memenuhi syarat 246G(iii), sehingga dapat
diintegralkan secara seragam.
      
\medskip 
      
{\bf (b)} Sekarang andaikan $A$ dapat diintegralkan secara seragam. Saya akan
mencari suatu himpunan $C\supseteq A$ yang kompak secara lemah.
      
\medskip 
      
\quad{\bf (i)} Untuk setiap $n\in\Bbb N$, pilih $E_n\in\Sigma$ dan $M_n\ge 0$
sedemikian sehingga $\mu E_n<\infty$ dan
$\int(|u|-M_n\chi E_n^{\ssbullet})^+\le 2^{-n}$ untuk setiap $u\in A$.
Tetapkan
      
\Centerline{$C=\{v:v\in L^1,\,|\int_Fv|\le M_n\mu(F\cap E_n) 
  +2^{-n}\Forall n\in\Bbb N,\,F\in\Sigma\}$,} 
      
\noindent dan perhatikan bahwa $A\subseteq C$. Memang, jika $u\in A$ dan
$F\in\Sigma$, maka
      
\Centerline{$|\int_Fu| 
\le\int_F(|u|-M_n\chi E_n^{\ssbullet})^++\int_FM_n\chi E_n^{\ssbullet} 
\le 2^{-n}+M_n\mu(F\cap E_n)$} 
      
\noindent untuk setiap $n$. Perhatikan pula bahwa $C$ terbatas dalam norma
$\|\,\|_1$, karena
      
\Centerline{$\|u\|_1\le 2\sup_{F\in\Sigma}|\int_Fu| 
\le 2\sup_{F\in\Sigma}(1+M_0\mu(F\cap E_0))\le 2(1+M_0\mu E_0)$} 
      
\noindent untuk setiap $u\in C$ (menggunakan 246F).
      
\medskip 
      
\quad{\bf (ii)} Karena teorema ini hendak dibuktikan untuk ruang ukur
sebarang $(X,\Sigma,\mu)$, kita tidak dapat menggunakan 243G untuk
mengidentifikasi dual $L^1$. Meskipun demikian, 247B di atas menunjukkan bahwa
243Gb `hampir' berlaku, dalam arti berikut: jika $h\in(L^1)^*$, terdapat
$v\in L^{\infty}$ sedemikian sehingga $h(u)=\int u\times v$ untuk setiap
$u\in C$. \Prf\ Tetapkan $G=\bigcup_{n\in\Bbb N}E_n\in\Sigma$, dan definisikan
$U\subseteq L^1$ seperti dalam 247A-247B. Menurut 247B, terdapat
$v\in L^{\infty}$ sedemikian sehingga $h(u)=\int u\times v$ untuk setiap
$u\in U$. Jika $u\in C$, nyatakan $u$ sebagai $f^{\ssbullet}$ dengan
$f:X\to\Bbb R$ terukur. Jika $F\in\Sigma$ dan $F\cap G=\emptyset$, maka
      
\Centerline{$|\int_Ff|=|\int_Fu|\le 2^{-n}+M_n\mu(F\cap E_n) 
=2^{-n}$} 
      
\noindent untuk setiap $n\in\Bbb N$, sehingga $\int_Ff=0$. Akibatnya,
$f=0$ hampir di mana-mana pada $X\setminus G$ (131Fc), jadi
$f\times\chi G\eae f$ dan $u=u\times\chi G^{\ssbullet}$; dengan kata lain,
$u\in U$, dan $h(u)=\int u\times v$, seperti yang diperlukan.
\Qed 
            
\medskip 
      
\quad{\bf (iii)} Dengan demikian kita dapat melanjutkan dengan suatu deskripsi
yang memadai, bukan tentang $(L^1(\mu))^*$ itu sendiri, melainkan tentang
tindakannya pada $C$.
      
Misalkan $\Cal F$ sebarang ultrafilter pada $L^1$ yang memuat $C$ (lihat 2A3R).
Untuk setiap $F\in\Sigma$, tetapkan
      
\Centerline{$\nu F=\lim_{u\to\Cal F}\int_Fu$;} 
      
\noindent karena
      
\Centerline{$\sup_{u\in C}|\int_Fu|\le\sup_{u\in C}\|u\|_1<\infty$,} 
      
\noindent nilai ini terdefinisi dengan baik dalam $\Bbb R$ (2A3S(e-ii)). Jika
$E$, $F$ adalah anggota $\Sigma$ yang saling lepas, maka
$\int_{E\cup F}u=\int_Eu+\int_Fu$ untuk setiap $u\in C$, sehingga
      
\Centerline{$\nu(E\cup F) 
=\lim_{u\to\Cal F}\int_{E\cup F}u 
=\lim_{u\to\Cal F}\int_{E}u+\lim_{u\to\Cal F}\int_{F}u 
=\nu E+\nu F$} 
      
\noindent (2A3Sf). Jadi $\nu:\Sigma\to\Bbb R$ aditif. Selanjutnya, fungsional
ini kontinu sejati terhadap $\mu$. \Prf\ Diberikan $\epsilon>0$, ambil
$n\in\Bbb N$ sedemikian sehingga $2^{-n}\le\bover12\epsilon$, tetapkan
$\delta=\epsilon/2(M_n+1)>0$, dan perhatikan bahwa
      
\Centerline{$|\nu F|\le\sup_{u\in C}|\int_Fu| 
\le 2^{-n}+M_n\mu(F\cap E_n)\le\epsilon$} 
      
\noindent apabila $\mu(F\cap E_n)\le\delta$.\ \Qed\ Menurut Teorema
Radon-Nikod\'ym (232E), terdapat $f_0\in{\eusm L}^1$ sedemikian sehingga
$\int_Ff_0=\nu F$ untuk setiap $F\in\Sigma$. Tetapkan
$u_0=f_0^{\ssbullet}\in L^1$. Jika $n\in\Bbb N$ dan $F\in\Sigma$, maka
      
\Centerline{$|\int_Fu_0| 
=|\nu F|\le\sup_{u\in C}|\int_Fu|\le 2^{-n}+M_n\mu(F\cap E_n)$,} 
      
\noindent sehingga $u_0\in C$.
      
\medskip 
      
\quad{\bf (iv)} Pokoknya ialah bahwa $\Cal F$ konvergen ke $u_0$.
\Prf\ Misalkan $h\in(L^1)^*$. Maka terdapat $v\in L^{\infty}$ sedemikian
sehingga $h(u)=\int u\times v$ untuk setiap $u\in C$. Nyatakan $v$ sebagai
$g^{\ssbullet}$, dengan $g:X\to\Bbb R$ suatu fungsi terbatas dan terukur
terhadap $\Sigma$. Ambil $\epsilon>0$. Pilih
$a_0\le a_1\le\ldots\le a_n$ sedemikian sehingga
$a_{i+1}-a_i\le\epsilon$ untuk setiap $i$, sementara $a_0\le g(x)<a_n$ untuk
setiap $x\in X$. Tetapkan $F_i=\{x:a_{i-1}\le g(x)<a_i\}$ untuk
$1\le i\le n$, serta tetapkan $\tilde g=\sum_{i=1}^na_i\chi F_i$ dan
$\tilde v=\tilde g^{\ssbullet}$; maka $\|\tilde v-v\|_{\infty}\le\epsilon$.
Kita mempunyai
      
$$\eqalign{\int u_0\times\tilde v 
&=\sum_{i=1}^na_i\int_{F_i}u_0 
=\sum_{i=1}^na_i\nu F_i\cr 
&=\sum_{i=1}^{n}a_i\lim_{u\to\Cal F}\int_{F_i}u 
=\lim_{u\to\Cal F}\sum_{i=1}^na_i\int_{F_i}u 
=\lim_{u\to\Cal F}\int u\times\tilde v.\cr}$$ 
      
\noindent Akibatnya,
      
$$\eqalign{\limsup_{u\to\Cal F}|\int u\times v-\int u_0\times v| 
&\le|\int u_0\times v-\int u_0\times \tilde v| 
+\sup_{u\in C}|\int u\times v-\int u\times\tilde v|\cr 
&\le\|u_0\|_1\|v-\tilde v\|_{\infty} 
+\sup_{u\in C}\|u\|_1\|v-\tilde v\|_{\infty}\cr 
&\le2\epsilon\sup_{u\in C}\|u\|_1.\cr}$$ 
      
\noindent Karena $\epsilon$ sebarang,
      
$$\eqalign{\limsup_{u\to\Cal F}|h(u)-h(u_0)| 
&=\limsup_{u\to\Cal F}|\int u\times v-\int u_0\times v| 
=0.\cr}$$ 
      
\noindent Karena $h$ sebarang, $u_0$ merupakan limit $\Cal F$ di $C$ dalam
topologi lemah pada $L^1$.\ \Qed
 
Karena $\Cal F$ sebarang, $C$ kompak secara lemah dalam $L^1$, dan bukti selesai.
}%end of proof of 247C 
      
\leader{247D}{Akibat} Misalkan $(X,\Sigma,\mu)$ dan $(Y,\Tau,\nu)$ sebarang
dua ruang ukur, dan $T:L^1(\mu)\to L^1(\nu)$ suatu operator linear kontinu.
Maka $T[A]$ merupakan subhimpunan $L^1(\nu)$ yang dapat diintegralkan secara
seragam setiap kali $A$ merupakan subhimpunan $L^1(\mu)$ yang dapat
diintegralkan secara seragam.
      
\proof{ Pokoknya, $T$ kontinu terhadap topologi lemah yang bersangkutan (2A5If).
Jika $A\subseteq L^1(\mu)$ dapat diintegralkan secara seragam, menurut 247C
terdapat himpunan $C\supseteq A$ yang kompak secara lemah. Karena merupakan
citra suatu himpunan kompak di bawah peta kontinu, $T[C]$ harus kompak secara lemah
(2A3N(b-ii)); jadi, menurut arah lain dalam 247C, $T[C]$ dan $T[A]$ dapat
diintegralkan secara seragam.
}%end of proof of 247D 
      
\cmmnt{ 
\leader{247E}{$L^1$ kompleks} Tidak ada kesulitan ataupun kejutan dalam
membuktikan 247C bagi $L^1_{\Bbb C}$. Jika kita mengikuti bukti yang sama,
semuanya tetap berlaku; tentu saja, kita harus ingat untuk mengubah konstanta
ketika menerapkan 246F, atau lebih tepatnya 246K, pada bagian (b-i) dari bukti.
} 
      
\exercises{ 
\leader{247X}{Latihan dasar $\pmb{>}$(a)} Misalkan $(X,\Sigma,\mu)$ sebarang
ruang ukur. Tunjukkan bahwa jika $A\subseteq L^1=L^1(\mu)$ relatif kompak dalam
topologi lemah, maka
$\{v:v\in L^1,\,|v|\le |u|$ untuk suatu $u\in A\}$ juga relatif kompak dalam
topologi lemah.
%247C 
      
\spheader 247Xb Misalkan $(X,\Sigma,\mu)$ suatu ruang ukur. Pada
$L^1=L^1(\mu)$, definisikan pseudometrik $\rho_F$, $\rho'_w$ untuk
$F\in\Sigma$, $w\in L^{\infty}(\mu)$ dengan menetapkan
$\rho_F(u,v)=|\int_Fu-\int_Fv|$, $\rho'_w(u,v)=|\int u\times w-\int 
v\times w|$ untuk $u$, $v\in L^1$. Tunjukkan bahwa pada sebarang subhimpunan
yang terbatas dalam norma $\|\,\|_1$ dari $L^1$, topologi yang didefinisikan oleh
$\{\rho_F:F\in\Sigma\}$ sama dengan topologi yang dibangkitkan oleh
$\{\rho'_w:w\in L^{\infty}\}$.
%247C 
      
\sqheader 247Xc Tunjukkan bahwa untuk sebarang himpunan $X$, suatu subhimpunan
dari $\ell^1=\ell^1(X)$ kompak dalam topologi lemah pada $\ell^1$ jika dan
hanya jika subhimpunan itu kompak dalam topologi norma pada $\ell^1$.
\Hint{246Xd.} 
%247C 
      
\spheader 247Xd Gunakan argumen (a-ii) dalam bukti 247C untuk menunjukkan
secara langsung bahwa jika $A\subseteq\ell^1(\Bbb N)$ kompak lemah, maka
$\inf_{n\in\Bbb N}|u_n(n)|=0$ bagi setiap barisan $\sequencen{u_n}$ dalam $A$.
%247C 
      
\spheader 247Xe Misalkan $(X,\Sigma,\mu)$ dan $(Y,\Tau,\nu)$ ruang ukur, dan
$T:L^2(\nu)\to L^1(\mu)$ sebarang operator linear terbatas. Tunjukkan bahwa
$\{Tu:u\in L^2(\nu),\,\|u\|_2\le 1\}$ dapat diintegralkan secara seragam dalam
$L^1(\mu)$. \Hint{gunakan 244K untuk melihat bahwa
$\{u:\|u\|_2\le 1\}$ kompak lemah dalam $L^2(\nu)$.}
%247C
      
\leader{247Y}{Latihan lanjutan (a)} 
%\spheader 247Ya 
Misalkan $(X,\Sigma,\mu)$ suatu ruang ukur. Ambil $1<p<\infty$ dan
$M\ge 0$, lalu tetapkan $A=\{u:u\in L^p=L^p(\mu),\,\|u\|_p\le M\}$. Tuliskan
$\frak S_A$ untuk topologi konvergensi dalam ukuran pada $A$, yakni topologi
subruang yang diinduksi oleh topologi konvergensi dalam ukuran pada $L^0(\mu)$.
Tunjukkan bahwa jika $h\in(L^p)^*$, maka $h\restr A$ kontinu terhadap
$\frak S_A$; jadi, jika $\frak T$ adalah topologi lemah pada $L^p$, topologi
subruang $\frak T_A$ termuat dalam $\frak S_A$.
%247C
      
\spheader 247Yb Misalkan $(X,\Sigma,\mu)$ suatu ruang ukur dan
$\sequencen{u_n}$ suatu barisan dalam $L^1=L^1(\mu)$ sedemikian sehingga
$\lim_{n\to\infty}\int_Fu_n$ terdefinisi untuk setiap $F\in\Sigma$. Tunjukkan
bahwa $\{u_n:n\in\Bbb N\}$ konvergen lemah. \Hint{246Yh.}   
%Find 
%an alternative argument relying on 2A5J and the result of 246Yj. (?)
%247C
}%end of exercises 
      
\endnotes{\Notesheader{247} Dalam 247D dan 247Xa saya berusaha menunjukkan
kekuatan identifikasi antara kekompakan lemah dan keterintegralan seragam.
Bahwa citra kontinu suatu himpunan yang kompak secara lemah harus kompak secara lemah merupakan
hal lumrah dalam analisis fungsional; bahwa selubung solid suatu himpunan yang
dapat diintegralkan secara seragam juga dapat diintegralkan secara seragam
langsung mengikuti definisi. Namun, saya tidak melihat argumen sederhana yang
menunjukkan bahwa citra kontinu suatu himpunan yang dapat diintegralkan secara
seragam harus dapat diintegralkan secara seragam, atau bahwa selubung solid
suatu himpunan yang kompak secara lemah harus relatif kompak dalam topologi lemah.
(Untuk pernyataan pertama memang ada cara lain; lihat 371Xf dalam jilid
berikutnya.)
      
Ada dua gagasan penting dalam bukti 247C. Gagasan pertama, dalam bagian (a-ii)
dari bukti, berupa manipulasi barisan secara cermat; inilah argumen yang
diperlukan untuk menunjukkan bahwa suatu subhimpunan dari $\ell^1$ yang kompak
secara lemah juga kompak dalam norma. (Mungkin akan membantu jika Anda menuliskan
penyelesaian 247Xd.) Himpunan $F_{n(k)}$ dan vektor $u_k$ dipilih untuk meniru
keadaan ketika terdapat suatu barisan dalam $\ell^1$ dengan $u_k(k)=1$ untuk
setiap $k$. Indeks $k(i)$ dipilih agar `tonjolan' itu bergerak cukup cepat,
sehingga $u_{k(j)}(k(i))$ sangat kecil apabila $i\ne j$. Akan tetapi, ini
berarti $\sum_{i=0}^{\infty}u_{k(j)}(k(i))$ (yang bersesuaian dengan
$\int_Gu_{k(j)}$ dalam bukti) selalu cukup besar, sedangkan
$\sum_{i=0}^{\infty}v(k(i))$ akan kecil bagi setiap calon titik klaster $v$
dari $\sequence{j}{u_{k(j)}}$. Saya menggunakan teknik serupa dalam \S246;
bandingkan 246Yg.
      
Pada bagian lain dari bukti 247C, strateginya lebih jelas. Anggota $L^1$
bersesuaian dengan fungsional kontinu sejati pada $\Sigma$; keterintegralan
seragam dari $C$ membuat himpunan fungsional yang bersesuaian `kontinu sejati
secara seragam', sehingga setiap fungsional batas juga kontinu sejati dan,
melalui Teorema Radon-Nikod\'ym, memberikan suatu anggota $L^1$. Argumen
aproksimasi langsung ((b-iv) dalam bukti, dan 247Xb) menunjukkan bahwa
$\lim_{u\in\Cal F}\int u\times w=\int v\times w$ untuk setiap
$w\in L^{\infty}$. Bagi ukuran $\mu$ yang dapat dilokalkan, ini akan
menyelesaikan bukti. Dalam kasus umum kita memerlukan satu langkah lagi, yang
di sini dilakukan dalam 247A-247B; suatu subhimpunan $L^1$ yang dapat
diintegralkan secara seragam secara efektif hidup pada bagian $\sigma$-hingga
dari ruang ukur, sehingga kita dapat mengabaikan bagian lainnya dan
mengandaikan bahwa ruang ukur tersebut dapat dilokalkan.
      
Syarat (ii)-(iv) dalam 246G memperlihatkan bahwa kekompakan lemah dalam $L^1$
dapat dibahas secara efektif dengan barisan; lihat juga 246Yh. Perlu dicatat
bahwa ini merupakan ciri umum kekompakan lemah dalam ruang Banach (2A5J).
Tentu saja, rumusan dengan barisan saling lepas dalam 246G merupakan ciri khas
$L^1$ -- maksud saya, meskipun ada hasil serupa yang berlaku di tempat lain
(lihat {\smc Fremlin 74}, bab.\ 8), gagasan-gagasannya tampak paling jelas dan
paling kuat ketika diterapkan pada $L^1$.
}%end of notes 
      
\frnewpage 
      
